\section{USAMO 1984}
        \begin{enumerate}
        	\item (\textit{USAMO 1984}) The product of two of the four roots of the quartic equation $x^4 - 18x^3 + kx^2+200x-1984=0$ is $-32$. Determine the value of $k$.


 

	\item (\textit{USAMO 1984}) The geometric mean of any set of $m$ non-negative numbers is the $m$-th root of their product.


 $\quad (\text{i})\quad$ For which positive integers $n$ is there a finite set $S_n$ of $n$ distinct positive integers such that the geometric mean of any subset of $S_n$ is an integer?


 $\quad (\text{ii})\quad$ Is there an infinite set $S$ of distinct positive integers such that the geometric mean of any finite subset of $S$ is an integer?


 

	\item (\textit{USAMO 1984}) $P$, $A$, $B$, $C$, and $D$ are five distinct points in space such that $\angle APB = \angle BPC = \angle CPD = \angle DPA = \theta$, where $\theta$ is a given acute angle. Determine the greatest and least values of $\angle APC + \angle BPD$.


 

	\item (\textit{USAMO 1984}) A difficult mathematical competition consisted of a Part I and a Part II with a combined total of $28$ problems. Each contestant solved $7$ problems altogether. For each pair of problems, there were exactly two contestants who solved both of them. Prove that there was a contestant who, in Part I, solved either no problems or at least four problems.


 

	\item (\textit{USAMO 1984}) $P(x)$ is a polynomial of degree $3n$ such that


 
 \begin{eqnarray*} P(0) = P(3) = \cdots &=& P(3n) = 2, \\ P(1) = P(4) = \cdots &=& P(3n-2) = 1, \\ P(2) = P(5) = \cdots &=& P(3n-1) = 0, \quad\text{ and }\\ && P(3n+1) = 730.\end{eqnarray*}


 Determine $n$.


 

\end{enumerate}